\subsection{标准分解}

在本节中，假设环A是交换环 \(k\) 上的（结合且有单位的）代数。对于 \(n \geqslant 0\)，我们用 \(\mathbf{B}_n\) 表示在 \(k\) 上的 \((n + 2)\) 个都等于 A 的模的张量积。我们将其视为 (A, A)-双模，赋予其左（相应地，右）A-模结构，该结构由张量积的第一个（相应地，最后一个）因子的左（相应地，右）A-模结构诱导。

对于 \(n \geqslant 1\)，我们定义双模同态 \(d_n^i\)（对于 \(0 \leqslant i \leqslant n\)）和 \(d_n\)，它们从 \(B_n\) 到 \(B_{n-1}\)，由公式：

\[
d _ {n} ^ {i} (x _ {0} \otimes \ldots \otimes x _ {n + 1}) = x _ {0} \otimes \ldots \otimes x _ {i} x _ {i + 1} \otimes \ldots \otimes x _ {n + 1}, \quad 0 \leqslant i \leqslant n,
\]

\[
d _ {n} = \sum_ {i = 0} ^ {n} (- 1) ^ {i} d _ {n} ^ {i}.
\]

显然，

\[
d _ {n - 1} ^ {i} \circ d _ {n} ^ {j} = d _ {n - 1} ^ {j - 1} \circ d _ {n} ^ {i} \quad \text { for } \quad i <   j
\]

并且因此

\[
d _ {n - 1} \circ d _ {n} = \sum_ {0 \leqslant i <   j \leqslant n} (- 1) ^ {i + j} d _ {n - 1} ^ {i} \circ d _ {n} ^ {j} + \sum_ {0 \leqslant j \leqslant i \leqslant n - 1} (- 1) ^ {i + j} d _ {n - 1} ^ {i} \circ d _ {n} ^ {j} = 0.
\]

因此，如果我们设 \(\mathbf{B}_n = 0\) 于 \(n < 0\)，且 \(d_n = 0\) 于 \(n \leqslant 0\)，则序列 \((\mathbf{B}_n, d_n)\) 定义了 (A, A)-双模的复形（X，第 43 页），将其记为 \(\mathbf{B}(\mathbf{A})\)。对于每个左 A-模 \(\mathbf{M}\)，我们用 \(\mathbf{B}(\mathbf{A}, \mathbf{M})\) 表示由 \(\mathbf{B}_n \otimes_{\mathbf{A}} \mathbf{M}\) 和 \(d_n \otimes 1_{\mathbf{M}}\) 构成的复形，换言之即张量积复形 \(\mathbf{B}(\mathbf{A}) \otimes_{\mathbf{A}} \mathbf{M}_{*}\)；它是一个左A-模复形。

我们定义一个A-线性映射 \(\varepsilon_{\mathbf{M}}\)，它从 \(\mathbf{B}_{0}(\mathbf{A},\mathbf{M}) = (\mathbf{A}\otimes_{k}\mathbf{A})\otimes_{\mathbf{A}}\mathbf{M}\) 到 \(\mathbf{M}\)，由公式 \(\varepsilon_{\mathbf{M}}(a\otimes b\otimes m) = abm\)（\(a,b\in \mathbf{A},m\in \mathbf{M}\)）给出。我们有 \(\varepsilon_{\mathbf{M}}\circ d_1 = 0\) ，使得分次同态 \(\overline{\varepsilon}_{\mathbf{M}}:\mathbf{B}(\mathbf{A},\mathbf{M})\to \mathbf{M}\) ，它在次数 0 上与 \(\varepsilon_{\mathbf{M}}\) 一致，是 A-模复形的一个态射。

命题 12. —— 映射 \(\overline{\varepsilon}_{\mathbf{M}}: \mathbf{B}(\mathbf{A}, \mathbf{M}) \to \mathbf{M}\) 是 \(k\) -模复形的一个同伦等价。特别地，复形 \(\mathbf{B}(\mathbf{A}, \mathbf{M})\) 在 \(k\) 上是分裂的，且 \((\mathbf{B}(\mathbf{A}, \mathbf{M}), \overline{\varepsilon}_{\mathbf{M}})\) 是A-模M的一个左分解。

对于 \(n \geqslant 0\) ，让我们定义一个 \(k\) -线性映射 \(s_n: \mathbf{B}_n \to \mathbf{B}_{n+1}\) 通过以下公式：

\[
s _ {n} \left(x _ {0} \otimes \dots \otimes x _ {n + 1}\right) = 1 \otimes x _ {0} \otimes \dots \otimes x _ {n + 1} \quad \text { for } \quad x _ {0}, \dots , x _ {n + 1} \in A.
\]

它是右A-模的同态，满足以下恒等式：

\[
d _ {n + 1} ^ {i} \circ s _ {n} = s _ {n - 1} \circ d _ {n} ^ {i - 1} \quad \mathrm{for} \quad n \geqslant 1, \quad 1 \leqslant i \leqslant n + 1,
\]

\[
d _ {n + 1} ^ {0} \circ s _ {n} = 1 _ {\mathrm{B} _ {n}} \quad \text { for } \quad n \geqslant 1,
\]

因此我们有

\[
d _ {n + 1} \circ s _ {n} + s _ {n - 1} \circ d _ {n} = 1 _ {\mathrm{B} _ {n}} \quad \text { for } \quad n \geqslant 1.\tag{16}
\]

此外，我们有

\[
d _ {1} \circ s _ {0} (x _ {0} \otimes x _ {1}) = x _ {0} \otimes x _ {1} - 1 \otimes x _ {0} x _ {1} \quad \text { for } \quad x _ {0}, x _ {1} \in A.\tag{17}
\]

我们用 \(\eta: \mathbf{A} \to \mathbf{A} \otimes_k \mathbf{A}\) 表示由 \(\eta(a) = 1 \otimes a\) 定义的映射，且用 \(\overline{\eta}: \mathbf{A} \to \mathbf{B}(\mathbf{A})\) 表示与 \(\eta\) 在次数0上一致的复形的态射。显然 \(\overline{\varepsilon}_{\mathbf{A}} \circ \overline{\eta} = 1_{\mathbf{A}}\) ；公式（16）和（17）表明 \(d \circ s + s \circ d = 1_{\mathrm{B(A)}} - \overline{\eta} \circ \overline{\varepsilon}_{\mathbf{A}}\) 。设 \(\overline{\eta}_{\mathrm{M}} = \overline{\eta} \otimes 1_{\mathrm{M}}\) , \(d_{\mathrm{M}} = d \otimes 1_{\mathrm{M}}\) 与 \(s_{\mathrm{M}} = s \otimes 1_{\mathrm{M}}\)，我们推出 \(\overline{\varepsilon}_{\mathrm{M}} \circ \overline{\eta}_{\mathrm{M}} = 1_{\mathrm{M}}\) 与 \(d_{\mathrm{M}} \circ s_{\mathrm{M}} + s_{\mathrm{M}} \circ d_{\mathrm{M}} = 1_{\mathrm{B(A,M)}} - \overline{\eta}_{\mathrm{M}} \circ \overline{\varepsilon}_{\mathrm{M}}\)。换言之（X，第33页，定义5）， \(\overline{\varepsilon}_{\mathrm{M}}\) 是 \(k\) -模复形的一个同伦等价。命题的其他断言立即成立。

定义3. --- 左分解(B(A, M), \(\overline{\varepsilon}_{\mathrm{M}}\) ) 称为A-模M的标准分解。

若A和M是投射（分别为自由、平坦）k-模，则标准分解B(A, M)是M的投射（分别为自由、平坦）分解。

